\begin{abstract}
\large \textbf{Abstract.}
\normalsize
We introduce PivotRL, a post-training method that achieves the in-domain accuracy and out-of-domain retention of end-to-end reinforcement learning while consuming only existing SFT trajectories and requiring no full-trajectory rollouts. Post-training for agentic capabilities currently forces a choice between SFT, which degrades out-of-domain performance, and end-to-end RL, which preserves it but demands repeated long-horizon environment interaction. PivotRL resolves this tension by identifying informative turns in expert trajectories—which we call \emph{pivots}—and performing on-policy RL only at those turns, scoring completions by whether they are aligned rather than whether they match the expert exactly. We prove that this approach minimally perturbs the reference policy while improving expert agreement, explaining why it avoids out-of-domain degradation. On Qwen3-30B-A3B, PivotRL improves $\tau^2$-Bench from 44.35 to 63.81, SWE-Bench Verified from 19.07 to 32.67, Terminal-Bench from 5.42 to 20.00, and BrowseComp from 2.50 to 11.30, while maintaining near-zero average out-of-domain change (+0.20) versus −9.48 for SFT on the same data.
\end{abstract}


% \begin{abstract}
%     Reinforcement learning (RL) is the most direct way to improve tool-using language models from execution feedback, but \emph{end-to-end token-level} RL is often impractical: tool-use episodes are long, rollouts are sequential, and standard objectives propagate a single return signal through thousands of tokens even though the environment only reacts at tool/turn boundaries.
% We mitigate this long-horizon bottleneck with two complementary changes.

% First, we formulate agentic tool use as an \emph{episodic Markov decision process at the interaction-turn granularity}, where each step is a complete assistant turn (typically a structured tool call or a short natural-language act) and advantages are defined over turns rather than **only** tokens.
% This allows us to align credit assignment with the interface that actually induces transitions and enabling reuse of cached interaction prefixes instead of repeatedly regenerating full episodes.

% Second, under this turn-level MDP we observe that advantages are typically \emph{sparse}: only a small subset of states exhibit large advantage gaps, while most states have nearly flat action values.
% This motivates \textbf{PivotRL}, which concentrates policy optimization on these high-leverage \emph{pivot states} and applies GRPO-style updates only to the pivot-turn span.
% By combining turn-level credit assignment with pivot-focused updates, PivotRL avoids the need to repeatedly execute full end-to-end sequential rollouts for every gradient step---a key practical barrier for agentic RL in the open source community, rivaling larger models like GPT-OSS-120B.

% On the $\tau^2$ Airline/Retail/Telecom benchmarks, PivotRL improves Qwen3-30B-A3B-2507-Thinking from $44.35$ to $62.16$ average success within $300$ RL steps. Further, the same approach is able to improve Qwen3-30B-A3B-2507-Thinking from $8.75$ success rate to $17.50$ success rate on TerminalBench-Core, using the Terminus harness.
% \end{abstract}
